\section{Introduction}

Large language models (LLMs) have demonstrated remarkable capabilities across a wide range of tasks. 
However, they are prone to hallucination, generating fluent but unsupported or incorrect responses~\cite{ji2023survey,Ji_2023,lin2022truthfulqameasuringmodelsmimic}, particularly under conditions of uncertainty or incomplete information. 
Such behavior poses significant challenges for real-world applications where reliability is critical.

Existing approaches to mitigate hallucination primarily focus on improving knowledge access or post-hoc verification. 
For example, retrieval-augmented generation (RAG)~\cite{rag} enhances factual grounding by incorporating external documents, while various detection methods attempt to identify hallucinated outputs after generation~\cite{lin2022truthfulqameasuringmodelsmimic}. 
Despite their effectiveness, these approaches do not explicitly address a fundamental issue: LLMs often fail to make appropriate decisions under uncertainty, such as when to answer, when to abstain, and how confidently to respond.

In this work, we argue that hallucination is not merely a generation error, but a consequence of suboptimal decision-making under uncertainty. 
From this perspective, mitigating hallucination requires learning a policy that balances correctness, confidence, and usefulness, rather than solely improving token-level prediction.

To this end, we propose a behavior-oriented reinforcement learning framework that explicitly models hallucination as a decision problem. 
Our approach optimizes model behavior along three key dimensions:
(i) whether to answer or abstain, based on the availability of reliable information; 
(ii) how to construct responses, ensuring grounded and informative outputs; and 
(iii) how confident to be, avoiding overconfident generation under uncertainty.

We design a unified reward function that integrates multiple components: 
behavior alignment, which encourages correct decision boundaries between answering and abstaining; 
uncertainty modeling via entropy-based overconfidence penalties; 
response quality shaping to promote completeness, actionability, and the use of supporting evidence; 
and length regularization to prevent degenerate short responses or excessively verbose outputs.

To optimize this objective, we adopt a two-stage training framework consisting of preference-based initialization using Direct Preference Optimization (DPO)~\cite{dpo}, followed by reinforcement learning with Group Relative Policy Optimization (GRPO)~\cite{GRPO}, enabling stable and fine-grained policy optimization.

We conduct extensive experiments to evaluate the proposed method. 
The results show that our approach reduces hallucination by over 65\%, consistently outperforming strong baselines such as GPT-4o. 
Furthermore, the learned behavior generalizes effectively to unseen domains, improves response quality across multiple evaluation dimensions, and preserves general knowledge capability without catastrophic forgetting~\cite{EWC}.

These findings demonstrate that hallucination mitigation can be effectively achieved by modeling it as a behavior learning problem under uncertainty, rather than solely improving knowledge access or generation mechanisms.